\chapter{总览}
\dirtree{%
.1 SQL 取数 = 对表中的「行」做去留、对「列」做计算，再把结果拼起来.
.2 0\quad 基础对象层（不是目的，是后面四种能力的对象）.
.3 表 / 行 / 列.
.4 要点：先想清楚「最后这张表长什么样」，再倒着拆步骤.
.3 执行顺序：FROM/JOIN $\to$ WHERE $\to$ GROUP BY $\to$ HAVING $\to$ 窗口函数 $\to$ SELECT $\to$ DISTINCT $\to$ ORDER BY.
.4 要点：窗口函数在 HAVING 之后才算，不能直接在 WHERE/HAVING 里用它的结果，要再套一层子查询.
.2 能力一\quad 简单筛选：要不要这一行 $\to$ WHERE.
.3 WHERE\quad 条件不满足就整行丢弃；只能用原始列，不能用聚合/窗口函数.
.3 HAVING\quad 对 GROUP BY 之后的结果再筛一次；可以用聚合函数.
.3 子查询当筛选条件\quad WHERE col IN (SELECT \ldots) $\approx$ 集合的交/差.
.2 能力二\quad 复杂筛选但不删行：给这一行打标签/变形 $\to$ SELECT.
.3 计算字段\quad 用已有列的表达式算出新列.
.3 CASE WHEN\quad 条件分支但不丢行：把「判断」变成新列的取值.
.4 配合聚合：SUM(CASE WHEN \ldots THEN x ELSE 0 END) 按条件统计.
.3 字符串/日期函数\quad CONCAT/SUBSTRING/DATE\_FORMAT/DATE\_ADD 等单值变换.
.2 能力三\quad 行与行之间的数据运算：横着拼成一行 $\to$ JOIN.
.3 INNER/LEFT/RIGHT JOIN\quad 按 key 把两张表的列拼到同一行（增加列）.
.3 自连接 self join\quad 同一张表的不同行互相对照.
.3 UNION / UNION ALL\quad 多个查询结果纵向叠在一起（增加行，不是 JOIN）.
.2 能力四\quad 一列数据内部、跨行的运算：不拼表也能看别的行 $\to$ 窗口函数.
.3 聚合函数 + GROUP BY\quad 把同列多行压成一个值（会减少行数）.
.3 窗口函数 函数() OVER (PARTITION BY \ldots ORDER BY \ldots).
.4 排名：ROW\_NUMBER / RANK / DENSE\_RANK / NTILE（分桶排名）.
.4 累计：SUM/AVG/COUNT OVER（配合 ROWS/RANGE 帧子句）.
.4 偏移：LAG / LEAD 取上下行同列的值（自连接的替代品）.
}
\chapter{能力三}
\section{自链接}
\begin{example}
用自连接做法，通过 \texttt{DATEDIFF} 限定严格前一天：

\begin{lstlisting}[language=SQL]
SELECT w1.id
FROM Weather AS w1
JOIN Weather AS w2
ON DATEDIFF(w1.recordDate, w2.recordDate) = 1
WHERE w1.Temperature > w2.Temperature;
\end{lstlisting}
\end{example}

\begin{note}
\texttt{LAG} 只取上一行，不管日期实际差了几天，若数据有缺失日期会误判。自连接配合 \texttt{DATEDIFF} 可以严格限定"恰好前一天"，更加健壮。
\end{note}
\begin{itemize}
    \item \texttt{DATEDIFF(日期1, 日期2)}：返回日期1 $-$ 日期2的天数差，如 \texttt{DATEDIFF('2015-01-02', '2015-01-01') = 1}
    \item \texttt{DATE\_ADD(日期, INTERVAL n DAY)}：在日期上加 $n$ 天，$n$ 为负数则为减，如 \texttt{DATE\_ADD('2015-01-02', INTERVAL -1 DAY) = '2015-01-01'}
\end{itemize}
\begin{dxtips}
 ON 后面就是普通的条件表达式，不一定要等号，大于小于、函数运算都可以，只要能描述两行之间的关联关系就行。
\end{dxtips}
\chapter{能力四}
\section{聚合}
\begin{dxtips}
    
\textbf{GROUP BY 的本质}

\texttt{GROUP BY} 内部做了两件事：
\begin{enumerate}
    \item \textbf{分组}：把相同键（如 \texttt{(student\_id, subject\_name)}）的行归到一堆
    \item \textbf{循环}：对每一堆执行聚合函数（\texttt{COUNT}、\texttt{SUM}、\texttt{AVG}……）
\end{enumerate}

这个「循环」由数据库引擎完成，无需手写。只需声明：
\begin{itemize}
    \item 按什么分——\texttt{GROUP BY} 的键
    \item 怎么聚合——\texttt{SELECT} 中的聚合函数
\end{itemize}

这也是 SQL 与命令式语言最大的区别：\textbf{描述你要什么，不描述怎么做}。

\end{dxtips}
\begin{dxtips}
    count的时候count主值
\end{dxtips}
\begin{solution}
\subsection{条件表达式}
\begin{lstlisting}[language=SQL]
SELECT sgn.user_id,
    ROUND(AVG(CASE WHEN action = 'confirmed' THEN 1 ELSE 0 END), 2)
        AS confirmation_rate
FROM signups sgn
LEFT JOIN confirmations c
    ON sgn.user_id = c.user_id
GROUP BY sgn.user_id
\end{lstlisting}

\begin{note}
用 \texttt{AVG(CASE WHEN ... THEN 1 ELSE 0 END)} 代替 \texttt{COUNT/COUNT}：
\begin{itemize}
    \item \texttt{confirmed} 的行计为 1，其余计为 0
    \item \texttt{AVG} 直接得到确认率，避免除以 0 的问题
    \item 没有任何记录的用户，\texttt{AVG} 返回 \texttt{NULL}，可用 \texttt{IFNULL(..., 0)} 处理
\end{itemize}
\end{note}
\begin{dxtips}
   对不对条件表达式就是再最后面能起到筛选的作用同时他不会删除行也就是同时能筛选符合条件的和不符合条件的在比值方面尤为突出，特别是加权平均
\end{dxtips}
\end{solution}
\section{窗口函数}
\begin{example}
查找比前一天温度更高的日期（力扣197）。

使用窗口函数 \texttt{LAG} 获取前一天温度，再在外层筛选：

\begin{lstlisting}[language=SQL]
SELECT id
FROM (
    SELECT id,
           Temperature,
           LAG(Temperature, 1) OVER (ORDER BY recordDate) AS prev_temp
    FROM Weather
) AS t
WHERE Temperature > prev_temp;
\end{lstlisting}

内层子查询产生的临时表结构如下：

\begin{center}
\begin{tabular}{cccc}
\hline
id & recordDate & Temperature & prev\_temp \\
\hline
1 & 2015-01-01 & 10 & NULL \\
2 & 2015-01-02 & 25 & 10 \\
3 & 2015-01-03 & 20 & 25 \\
4 & 2015-01-04 & 30 & 20 \\
\hline
\end{tabular}
\end{center}

外层 \texttt{WHERE Temperature > prev\_temp} 筛选出 id = 2, 4。
\end{example}

\begin{note}
\texttt{LAG(字段, 偏移量) OVER (ORDER BY 排序字段)} 取当前行往前数第 $n$ 行的值，无上一行时自动返回 \texttt{NULL}。窗口函数不能直接在 \texttt{WHERE} 中引用，必须套一层子查询后才能筛选。
\end{note}
\begin{dxtips}
  本质上窗口函数还是把行之间的问题转化为列之间的问题  
\end{dxtips}
    
\begin{dxtips}
    由于窗口函数只能放进select和order里面如果最后不想看见他就只能再套一层子查询
\end{dxtips}
\begin{itemize}
    \item \texttt{LAG(字段, n) OVER (ORDER BY ...)}：从当前行向前取第 $n$ 行的值，当前行"滞后"于目标行
    \item \texttt{LEAD(字段, n) OVER (ORDER BY ...)}：从当前行向后取第 $n$ 行的值，当前行"领先"于目标行
\end{itemize}
\begin{dxtips}
    因为where 比 select 先执行所以外层因为select里面新出的东西不能进where
\end{dxtips}